\documentclass[letterpaper,10pt]{article}

\usepackage{latexsym}
\usepackage[empty]{fullpage}
\usepackage{titlesec}
\usepackage{marvosym}
\usepackage[usenames,dvipsnames]{color}
\usepackage{verbatim}
\usepackage{enumitem}
\usepackage[hidelinks]{hyperref}
\usepackage{fancyhdr}
\usepackage[brazilian]{babel}
\usepackage{tabularx}
\input{glyphtounicode}

\usepackage[sfdefault]{roboto}

\pagestyle{fancy}
\fancyhf{}
\fancyfoot{}
\renewcommand{\headrulewidth}{0pt}
\renewcommand{\footrulewidth}{0pt}
\setlength{\footskip}{5pt}

\addtolength{\oddsidemargin}{-0.5in}
\addtolength{\evensidemargin}{-0.5in}
\addtolength{\textwidth}{1in}
\addtolength{\topmargin}{-.5in}
\addtolength{\textheight}{1.0in}

\urlstyle{same}
\raggedbottom
\raggedright
\setlength{\tabcolsep}{0in}

\titleformat{\section}{\Large\bfseries\scshape\raggedright}{}{0em}{}[\titlerule]

\pdfgentounicode=1

\newcommand{\resumeItem}[1]{\item\small{#1}}
\newcommand{\resumeSubheading}[4]{
\vspace{-1pt}\item
  \begin{tabular*}{0.97\textwidth}[t]{l@{\extracolsep{\fill}}r}
    \textbf{#1} & #2 \\
    \textit{#3} & \textit{#4} \\
  \end{tabular*}\vspace{-7pt}
}
\renewcommand\labelitemii{$\vcenter{\hbox{\tiny$\bullet$}}$}
\newcommand{\resumeSubHeadingList}{\begin{itemize}[leftmargin=0.15in, label={}]}
\newcommand{\resumeSubHeadingListEnd}{\end{itemize}}

\begin{document}

\begin{center}
  \textbf{\Huge Gabriel Frigo} \\
  \small Santo André, SP $|$ +55 11 95570-2000 $|$ \href{mailto:gabriel.frigo4@gmail.com}{gabriel.frigo4@gmail.com} $|$
  \href{https://linkedin.com/in/gabriel-frigo-b6727b275/}{linkedin.com/in/gabriel-frigo} \\
  \href{https://github.com/GabrielFrigo4}{github.com/GabrielFrigo4} $|$
  \href{https://gabrielfrigo.dev.br}{gabrielfrigo.dev.br}
\end{center}

\section{Resumo Profissional}
Estudante de Ciência da Computação na UFABC com sólida formação algorítmica e matemática (Medalhista de Ouro OBMEP, 13º no Brasil; competidor ICPC/SBC), conectando Pesquisa Operacional e rigor acadêmico à engenharia de produção em larga escala. Perfil autodidata focado em \textbf{Inteligência Artificial, Engenharia Backend e Sistemas Escaláveis}, unindo o ecossistema moderno de nuvem e microsserviços (Go, Python, SQL, Clojure, Docker/Podman, OpenTofu IaC) aos fundamentos de sistemas e ecossistemas Unix (C/C++, FreeBSD, POSIX). Praticante avançado da cultura \textit{AI-First} e desenvolvimento agêntico (\textbf{Antigravity IDE/CLI, Claude, OpenAI APIs}) para prototipagem rápida e resolução de problemas de alta complexidade.

\section{Competências Técnicas}
\resumeSubHeadingList
  \resumeItem{\textbf{Linguagens de Programação}: Go, Python, SQL, Clojure, Java, C, C++, Rust, C\#, Zig, Bash, Lua}
  \resumeItem{\textbf{IA, Otimização \& Agentes}: Google OR-Tools, APIs de LLMs (OpenAI, Claude, Gemini), Antigravity IDE/CLI, Cursor, Automação Inteligente, Pandas, NumPy}
  \resumeItem{\textbf{Backend, Cloud \& IaC}: OpenTofu/Terraform (IaC), Magalu Cloud, Oracle Cloud (OCI), RESTful APIs, PostgreSQL, SQLite, PocketBase, Podman, Docker, Linux, FreeBSD, Caddy, Nginx}
  \resumeItem{\textbf{Engenharia \& Ferramentas}: Git/Got (Game of Trees), GitHub Actions (CI/CD), POSIX Make, CGo / Interoperabilidade, Pesquisa Operacional (VRPTW), Otimização Combinatória, Estruturas de Dados Avançadas, Microsserviços, Concorrência, Design Patterns}
  \resumeItem{\textbf{Competências Pessoais}: Raciocínio analítico afiado, autonomia com foco em entrega, adaptabilidade técnica rápida, comunicação transparente e colaboração}
  \resumeItem{\textbf{Idiomas}: Português (Nativo); Inglês (Avançado / Intermediário -- Certificação EFSET)}
\resumeSubHeadingListEnd

\section{Formação Acadêmica}
\resumeSubHeadingList
  \resumeSubheading
    {Universidade Federal do ABC (UFABC)}{Santo André, SP}
    {Bacharelado em Ciência da Computação}{Jun/2023 -- Jun/2028}
  \resumeSubheading
    {Universidade Federal do ABC (UFABC)}{Santo André, SP}
    {Bacharelado em Ciência e Tecnologia}{Jun/2023 -- Dez/2027}
\resumeSubHeadingListEnd

\section{Experiência e Pesquisa em Engenharia}
\resumeSubHeadingList

  \resumeSubheading
    {\href{https://github.com/GabrielFrigo4/IC_Networks_Flow}{Otimização e Fluxo em Redes: Teoria e Algoritmos}}{Out/2025 -- Set/2026}
    {Iniciação Científica (CNPq / UFABC)}{Santo André, SP}
    \resumeSubHeadingList
      \resumeItem{\textbullet\ Pesquisa em Teoria dos Grafos, Algoritmos de Fluxo Máximo e Otimização Combinatória voltados para eficiência e escalabilidade algorítmica.}
      \resumeItem{\textbullet\ Implementação de algoritmos e estruturas de dados de alta performance em C++, com análise formal de complexidade assintótica e benchmarks.}
    \resumeSubHeadingListEnd

  \resumeSubheading
    {Programação Competitiva GRUB (ICPC / Maratona SBC)}{Jan/2024 -- Presente}
    {Membro e Competidor}{UFABC, SP}
    \resumeSubHeadingList
      \resumeItem{\textbullet\ Resolução de problemas algorítmicos complexos sob limites rígidos de tempo de execução e memória utilizando C/C++ e Python.}
      \resumeItem{\textbullet\ Aplicação de programação dinâmica, grafos, geometria computacional e estruturas de dados avançadas em competições oficiais.}
    \resumeSubHeadingListEnd

  \resumeSubheading
    {Green Team Hacker Club (GTHC)}{Dez/2023 -- Presente}
    {Analista de Projetos e Instrutor Técnico}{Santo André, SP}
    \resumeSubHeadingList
      \resumeItem{\textbullet\ Mentoria e instrução técnica para 10+ estudantes em arquitetura de sistemas, engenharia de software e versionamento com Git.}
      \resumeItem{\textbullet\ Ensino de C/C++, Rust e fundamentos de baixo nível aplicados à construção de software eficiente e performático.}
    \resumeSubHeadingListEnd

  \resumeSubheading
    {\href{https://github.com/GabrielFrigo4/TT_MiniSumo_Arruela}{Robótica Competitiva Tamandutech}}{Jun/2023 -- Abr/2026}
    {Desenvolvedor de Software}{UFABC, SP}
    \resumeSubHeadingList
      \resumeItem{\textbullet\ Desenvolvimento de lógica de controle em tempo real e algoritmos autônomos para tomada de decisão em robôs de competição.}
      \resumeItem{\textbullet\ Repositório adicional: \href{https://github.com/GabrielFrigo4/TTCC}{TTCC}.}
    \resumeSubHeadingListEnd

  \resumeSubheading
    {\href{https://github.com/GabrielFrigo4/PDPD_Ramsey}{Iniciação Científica PICME (IMPA / OBMEP)}}{Mar/2023 -- Jan/2026}
    {Pesquisador Acadêmico}{UFABC, SP}
    \resumeSubHeadingList
      \resumeItem{\textbullet\ Pesquisa em Teoria de Ramsey e Combinatória Infinita, desenvolvendo raciocínio matemático rigoroso aplicado a estruturas de redes e grafos.}
    \resumeSubHeadingListEnd

\resumeSubHeadingListEnd

\section{Projetos de Software \& IA}
\resumeSubHeadingList

  \resumeSubheading
    {\href{https://github.com/GabrielFrigo4/optilaser}{OptiLaser (Otimização Logística)}}{Jul/2026 -- Presente}
    {Pesquisa Operacional, Backend \& Cloud IaC}{Go/CGo, OR-Tools, PocketBase, OpenTofu, Magalu Cloud}
    \resumeSubHeadingList
      \resumeItem{\textbullet\ Resolução do problema de roteamento e alocação médica (VRPTW) em $<3$s com \textbf{Google OR-Tools}.}
      \resumeItem{\textbullet\ Backend em Go e CGo multiplataforma (\textbf{FreeBSD/Linux}) com banco PocketBase embarcado.}
      \resumeItem{\textbullet\ Deploy via \textbf{Podman} com \textbf{Rocky Linux} provisionado por \textbf{OpenTofu} na \textbf{Magalu Cloud}.}
    \resumeSubHeadingListEnd

  \resumeSubheading
    {\href{https://github.com/GabrielFrigo4/unix-sock}{Servidor Web HTTP Concorrente}}{Fev/2026 -- Mai/2026}
    {Redes, Sistemas Concorrentes e Performance}{C, POSIX/BSD Sockets, Concorrência, Caddy}
    \resumeSubHeadingList
      \resumeItem{\textbullet\ Servidor HTTP implementado do zero em C com chamadas diretas à API de POSIX Sockets e suporte multiplataforma (Linux/FreeBSD).}
      \resumeItem{\textbullet\ Arquitetura concorrente por multiprocessamento para tratamento simultâneo de requisições web com alta vazão e baixa latência.}
      \resumeItem{\textbullet\ Deploy em produção na nuvem (OCI) com Caddy proxy; Projeto live: \href{https://game.gabrielfrigo.dev.br/}{game.gabrielfrigo.dev.br}}
    \resumeSubHeadingListEnd

  \resumeSubheading
    {Mentoria em Algoritmos (Projeto Extensão OBI / UFABC)}{Jan/2026 -- Presente}
    {Instrutor de Programação}{UFABC, SP}
    \resumeSubHeadingList
      \resumeItem{\textbullet\ Treinamento de 10+ alunos em lógica, estruturas de dados e resolução de problemas computacionais para a Olimpíada Brasileira de Informática.}
    \resumeSubHeadingListEnd

\resumeSubHeadingListEnd

\section{Prêmios e Conquistas Acadêmicas}
\resumeSubHeadingList
  \resumeItem{\textbf{Medalha de Ouro (13ª Posição Nacional, 2022)}, Prata (2018, 2019) e Bronze (2017, 2021) -- OBMEP}
  \resumeItem{\textbf{14ª Posição (2026) e 15ª Posição (2025)} -- Maratona Paulista de Programação}
  \resumeItem{\textbf{107ª Posição (2025) e 135ª Posição (2024) na Fase Regional} -- Maratona SBC de Programação / ICPC}
  \resumeItem{\textbf{Menção Honrosa (2024)} -- Olimpíada Brasileira de Informática (OBI)}
  \resumeItem{\textbf{Medalha de Prata (2023) e Menção Honrosa (2025)} -- Competição de Lógica (CELLM)}
  \resumeItem{\textbf{Medalha de Bronze (2021)} -- Olimpíada Paulista de Matemática (OPM)}
\resumeSubHeadingListEnd

\section{Certificações}
\resumeSubHeadingList
  \resumeItem{\textbullet\ \href{https://efset.org/cert/9uJjv8}{EFSET English Certificate (Intermediário / Avançado)}}
\resumeSubHeadingListEnd

\end{document}
